\documentclass[10pt]{article}
\usepackage[margin=0.76in]{geometry}
\usepackage{amsmath,amssymb,amsthm}
\usepackage[hidelinks]{hyperref}
\setlength{\parindent}{0pt}
\setlength{\parskip}{0.42em}

\newtheorem{theorem}{Theorem}
\newtheorem{remark}[theorem]{Remark}
\newcommand{\Npos}{\mathbb N_{+}}
\newcommand{\Torus}{\mathbb T}

\begin{document}

\begin{center}
{\LARGE Positive Cosine Powers Disprove the}\par\vspace{0.15em}
{\LARGE Finite-$L^p$ Form of Conjecture 1}\par\vspace{0.4em}
{\large A counterexample to arXiv:1606.09092}\par\vspace{0.2em}
August 9, 2026
\end{center}

\begin{center}
\fbox{\parbox{0.91\textwidth}{\textbf{Status.}
Candidate counterexample, likely valid. It disproves the conjecture exactly as
written by exposing the role of the omitted zeroth power in finite $L^p$; the
corrected version with $0$ imposed remains open in this packet.}}
\end{center}

\section*{The conjecture}

Jaming and Simon~\cite{JamingSimon} define a subset
$\Lambda\subset\mathbb N=\{0,1,2,\ldots\}$ to be a
$[-1,1]$-M\"untz--Sz\'asz sequence only if $0\in\Lambda$ and the reciprocal
sums over its positive even and odd elements both diverge. Their Conjecture 1
(Conjecture 4.6 in the published numbering) says that, for
$p\in[1,\infty]$ and $\theta_1-\theta_2\notin\mathbb Q$, totality of
\[
 \{\cos^\lambda 2\pi(t-\theta_1):\lambda\in\Lambda\}
 \ \cup\
 \{\cos^{\lambda'}2\pi(t-\theta_2):\lambda'\in\Lambda'\}
 \tag{1}
\]
in $X_p(0,1)$ forces both $\Lambda$ and $\Lambda'$ to satisfy that definition.

For finite $p$, a single point carries no $L^p$ information, and the condition
$0\in\Lambda$ is not necessary. This gives a counterexample.

\begin{theorem}[counterexample]\label{thm:counterexample}
Let $1\le p<\infty$ and let $\theta_1,\theta_2\in[0,1)$ satisfy
$\theta_1-\theta_2\notin\mathbb Q$. Set
\[
                  \Lambda=\Lambda'=\Npos:=\{1,2,3,\ldots\}.
\]
Then the system in \emph{(1)} is total in $L^p(0,1)$, although neither
$\Lambda$ nor $\Lambda'$ is a $[-1,1]$-M\"untz--Sz\'asz sequence under the
definition of~\cite{JamingSimon}.
\end{theorem}

\begin{proof}
Identify $[0,1)$ with the circle $\Torus$. For $\theta\in\Torus$, let
\[
 E_\theta^p=\{f\in L^p(\Torus):f(\theta+t)=f(\theta-t)\text{ a.e.}\}
\]
be the closed reflection-even subspace. Every positive power of
$\cos 2\pi(t-\theta)$ belongs to $E_\theta^p$.

The missing constant is nevertheless in their $L^p$-closed span. Indeed,
\[
 q_N(t)=1-\bigl(1-\cos^2 2\pi(t-\theta)\bigr)^N
       =\sum_{j=1}^{N}(-1)^{j+1}\binom Nj
          \cos^{2j}2\pi(t-\theta).
 \tag{2}
\]
We have $0\le q_N\le1$, and $q_N(t)\to1$ except at the two zeros of the cosine.
Dominated convergence gives $q_N\to1$ in $L^p(\Torus)$.

It follows that the positive powers and all nonnegative powers have the same
closed span. That span is exactly $E_\theta^p$: one inclusion is reflection
symmetry, while the reverse inclusion follows because the Chebyshev identity
\[
 \cos 2\pi k(t-\theta)=T_k\!\left(\cos 2\pi(t-\theta)\right)
\]
puts every reflection-even trigonometric polynomial in the closed span, and
such polynomials are dense in $E_\theta^p$. Hence the closed spans of the two
families in (1) are respectively $E_{\theta_1}^p$ and $E_{\theta_2}^p$.

We now show that $E_{\theta_1}^p+E_{\theta_2}^p$ is dense. Let
$g\in L^{p'}(\Torus)$ annihilate both subspaces, where $p'$ is the conjugate
index (with $p'=\infty$ when $p=1$). Since $g\in L^1(\Torus)$, its Fourier
coefficients are defined. Annihilation of
$\cos 2\pi k(t-\theta_j)\in E_{\theta_j}^p$ gives, for $j=1,2$ and $k\ge1$,
\[
 e^{-2\pi i k\theta_j}\widehat g(-k)
   +e^{2\pi i k\theta_j}\widehat g(k)=0.
 \tag{3}
\]
Eliminating $\widehat g(-k)$ between the two equations yields
\[
 \left(e^{4\pi i k\theta_1}-e^{4\pi i k\theta_2}\right)\widehat g(k)=0.
\]
The factor in parentheses never vanishes for $k\ne0$, because
$\theta_1-\theta_2$ is irrational. Thus $\widehat g(k)=\widehat g(-k)=0$
for every $k\ge1$. The constant belongs to both even subspaces, so
$\widehat g(0)=0$ as well. Uniqueness of Fourier coefficients in $L^1$ gives
$g=0$. Hahn--Banach now proves totality of (1).

Finally, $0\notin\Npos$. Therefore neither chosen exponent set meets the
source's definition of a $[-1,1]$-M\"untz--Sz\'asz sequence, completing the
counterexample.
\end{proof}

\section*{Exact scope}

The counterexample works for every finite $p$, including $p=1$, and its two
reciprocal parity sums do diverge; failure of the source condition is solely
the missing exponent $0$. This is the familiar distinction between uniform
and $L^p$ approximation: (2) cannot converge uniformly, because every $q_N$
vanishes where the cosine vanishes. Thus no claim is made against the
$C([0,1])$ endpoint or against the corrected finite-$p$ conjecture in which
$0\in\Lambda\cap\Lambda'$ is imposed explicitly.

Bounded searches through the source title, its exact conjecture, and the two
shifted cosine-power terminology found no later paper recording this
counterexample. It should nevertheless be viewed as a correction of the
statement rather than a resolution of the intended small-divisor problem.

{\small
\begin{thebibliography}{9}
\bibitem{JamingSimon}
P.~Jaming and I.~Simon,
\emph{M\"untz--Sz\'asz type theorems for the density of the span of powers of
functions}, Bull. Sci. Math. 166 (2021), 102933; arXiv:1606.09092.
\end{thebibliography}
}

\end{document}

